\section{The problem: a distant legislature and a contested record}
\label{sec:context}

This section describes the setting the system was built for. Every figure is taken from the source cited beside it; all web sources were read on 16 November 2026.

\subsection{The distance between citizens and their deputies}

In June 2026, 67\% of Brazilian voters did not remember whom they had voted for federal deputy four years earlier, and 68\% could not name a single sitting deputy. Only 7\% had forgotten their vote for president. Six of the 513 deputies were named by at least 1\% of respondents~\citep{folha2026votememory}. The pattern is old: in the 2014 national election study, more than 40\% could not recall their vote for deputy a few weeks after casting it~\citep{fapesp2015eseb}.

Distance comes with distrust. In September 2026, 45\% of voters said they do not trust the National Congress and 45\% that they do not trust political parties~\citep{datafolha2026trust}; on a 0--100 index of social trust the Congress scored 34 and the parties 30, the two worst results of the survey~\citep{ipsos2026ics}. Support for democracy in Brazil stood at 41\%, the second lowest among the Latin American countries surveyed~\citep{latinobarometro2026}, and 37\% of adults were satisfied with how democracy works~\citep{pew2026brazil}. A quarter of voters declared no interest in politics at all~\citep{datafolha2026interest}.

Part of the distance is institutional. The Chamber has 513 members elected by open-list proportional representation in state-wide districts~\citep{tse2026systems}, a system that, in one influential account, produces weak parties and individualistic politicians with little accountability to citizens~\citep{ames2001deadlock}, and in which few voters can follow the path from their vote to a seat~\citep{nicolau2017representantes}. Twenty parties won seats in 2026, down from thirty in 2018~\citep{folha2026fragmentation}; the effective number of parties fell from 16.7 to 8.8 over the same period~\citep{tejadas2026nep}, from a level that had made the Chamber the most fragmented lower house in the world~\citep{poder3602023fragmentation}. Presidents govern through broad, shifting coalitions~\citep{abranches1988presidencialismo}, so the question of who supports the government has a different answer in each legislature. The largest party in the Chamber elected in 2026 holds 121 seats, under a quarter of the house~\citep{cnn2026seats}.

Part of it is volume. Between 1988 and 2017, 56,804 ordinary bills were filed in the Chamber, of which 4.63\% were approved~\citep{cerdeira2018congresso}. Much of what is enacted is symbolic: of the 164 ordinary laws of 2010, 43 were tributes or commemorative dates~\citep{queiroz2010producao}. A citizen who counts bills filed measures activity, not results.

Part of it is opacity that the courts themselves have had to correct. Parliamentary amendments to the federal budget amount to about R\$\,61 billion in 2026~\citep{senado2026budget}. In December 2022 the Supreme Court struck down one category of them, the so-called secret budget, because the amendments were anonymous and so violated the constitutional principle of publicity~\citep{stf2022secretbudget}.

And part of it is language. In the most recent national survey, 29\% of Brazilians aged 15 to 64 were functionally illiterate and only 10\% were proficient readers~\citep{inaf2024}. A federal law of November 2025 created a national policy of \emph{simple language} (\emph{linguagem simples}) and requires every public body, in all branches and levels of government, to use it with citizens. The law defines simple language as the set of techniques for conveying information clearly and objectively, so that the words, structure and layout of a message let citizens easily find the information, understand it and use it~\citep{brasil2025lei15263}. We use the term in that sense throughout: not a register or a style, but the removal of abstractions that stand between a citizen and the information. Legislative records, written in procedural terms for specialists, are far from that standard. In the vocabulary of \citet{michener2013identifying}, they are \emph{visible} but not \emph{inferable}: the public can see them and still cannot draw accurate conclusions from them.

\subsection{Contested elections and the disinformation that surrounds them}

The distance would matter less if citizens trusted those who report on the legislature. Brazil's last two presidential elections show why they often do not.

In 2022, Luiz In\'acio Lula da Silva defeated the incumbent, Jair Bolsonaro, with 50.90\% of the valid votes, the narrowest margin since direct presidential elections resumed in 1989~\citep{tse2022result,g12022margin}. In a panel study around that election, about half of respondents rated as accurate the false statements that votes are counted in a secret room and that the voting machines cannot be audited~\citep{montalverne2024nexus}. A week after the inauguration, on 8 January 2023, rioters invaded the seats of the three branches of government~\citep{stf8jan}. Bolsonaro was later declared ineligible for attacking the voting system before foreign ambassadors~\citep{tse2023ineligible} and, in September 2025, sentenced to 27 years and three months in prison for crimes that include an attempted coup~\citep{stf2025conviction}. In October 2026 Lula won again, with 51.79\% against Fl\'avio Bolsonaro~\citep{tse2026runoff}, who had led the first round~\citep{g12026rounds}; the defeated candidate refused to concede and alleged fraud without presenting evidence~\citep{reuters2026concede,cnn2026letter}, while international observers reported no doubts about the integrity of the results~\citep{oas2026brazil}. We take no position on any candidate. What matters here is that, twice in four years, close to half the electorate saw leaders it supports cast doubt on an official record, and had no practical way to check it.

Disinformation filled that space through the channels people actually use. In 2018, images that fact-checkers had confirmed as false were fewer than 1\% of the images shared in monitored WhatsApp political groups but reached 44\% of the groups~\citep{resende2019whatsapp}, and companies were reported to be paying for bulk messaging against one party~\citep{camposmello2018empresarios}. In 2022, having joined a political WhatsApp group was the strongest predictor of believing electoral falsehoods~\citep{rossini2023electoral}, and public groups were used to organise the attack of 8 January~\citep{ozawa2024capitol}. In 2026 the content became synthetic: one fact-checking project confirmed 1,423 political items made with generative AI in the campaign, about 21 a day, of which 768 carried no label~\citep{lupa2026vigia}, despite the electoral court's rules on synthetic content~\citep{tse2024res23732,tse2026res23755}. The bill that would regulate platforms has waited for a floor vote since 2023~\citep{brasil2020pl2630}.

Corrections struggle against this. Fact-checks do reduce false beliefs on average~\citep{walter2020factchecking,porter2021global}, but falsehood spreads faster than truth~\citep{vosoughi2018spread}, 40.7\% of the shares of debunked images in Brazilian WhatsApp groups happened after the debunking~\citep{reis2020debunked}, and an experiment during the 2018 Brazilian election found corrections ``ineffective at reducing misinformation''~\citep{batistapereira2022fake}. Underneath is a problem of trust: in 2026 trust in news in Brazil fell to 36\%, a twelve-year low, while 53\% of internet users got news from social media each week and 46\% from WhatsApp~\citep{carro2026brazil}. A correction asks its reader to believe the corrector, and that request often fails before the evidence is seen.

\subsection{The gap}

The Chamber is not secretive. It publishes its records as structured open data, and Brazil's access-to-information law makes publicity the rule and secrecy the exception~\citep{brasil2011lai}. The gap is between raw publication and public understanding, and it is filled by intermediaries: newspapers, rankings, influencers and message groups. Each has a point of view, or is assumed to have one, and none can be checked by its audience.

\citet{oneill2002question} argued two decades ago that transparency alone does not build trust, and that trust has a basis only when people are offered not ``a flood of uncheckable information, but rather information whose accuracy they can check and assess for themselves''. The system described next is an attempt to build an intermediary of that second kind: one that brings the legislature closer in simple language, and that does not ask to be believed.
